% ── Sección: Uso ──────────────────────────────────────────────────────────────
\section{Uso}

% ── Flujo de trabajo ──────────────────────────────────────────────────────────
\begin{frame}[fragile]{Flujo de trabajo completo}
  \begin{center}
    \begin{tikzpicture}[
      node distance=2.6cm,
      wfnode/.style={rectangle, draw, rounded corners=5pt,
                     minimum width=2.8cm, minimum height=1cm,
                     fill=pptexBlue!15, font=\small\bfseries},
      lbl/.style={font=\tiny\ttfamily, text=gray, below=2pt},
    ]
      \node[wfnode] (new)  {new};
      \node[wfnode] (edit) [right of=new]    {editar};
      \node[wfnode] (comp) [right of=edit]   {compile};
      \node[wfnode] (pdf)  [right of=comp]   {PDF};

      \draw[->, thick] (new)  -- (edit);
      \draw[->, thick] (edit) -- (comp);
      \draw[->, thick] (comp) -- (pdf);
      \draw[->, thick, dashed, bend right=35] (comp) to node[below, font=\tiny]{watch} (edit);

      \node[lbl] at (new.south)  {pptex new};
      \node[lbl] at (comp.south) {pptex compile};
    \end{tikzpicture}
  \end{center}

  \vspace{0.5em}
  El modo \texttt{watch} recompila automaticamente cada vez que guardas el archivo.
\end{frame}

% ── Proyectos ─────────────────────────────────────────────────────────────────
\begin{frame}[fragile]{Crear y organizar proyectos}
  \textbf{Diapositivas:}
  \begin{lstlisting}[language=bash]
./pptex new slides congreso
./pptex new slides curso_estadistica
  \end{lstlisting}

  \textbf{Informes:}
  \begin{lstlisting}[language=bash]
./pptex new report tesis
./pptex new report informe_final
  \end{lstlisting}

  \textbf{Ver todos los proyectos:}
  \begin{lstlisting}[language=bash]
./pptex list
  \end{lstlisting}

  Cada proyecto queda en su propia subcarpeta:
  \texttt{slides/congreso/}, \texttt{reports/tesis/}, etc.
\end{frame}

% ── Compilar ──────────────────────────────────────────────────────────────────
\begin{frame}[fragile]{Compilar}

  \textbf{Compilacion basica (XeLaTeX):}
  \begin{lstlisting}[language=bash]
./pptex compile slides congreso
  \end{lstlisting}

  \textbf{Con otro motor y bibliografia:}
  \begin{lstlisting}[language=bash]
./pptex compile reports tesis --engine lualatex --biber
  \end{lstlisting}

  \textbf{Modo watch — recompila al guardar:}
  \begin{lstlisting}[language=bash]
./pptex watch slides congreso
  \end{lstlisting}

  \textbf{Shell interactivo dentro del contenedor:}
  \begin{lstlisting}[language=bash]
./pptex shell
  \end{lstlisting}
\end{frame}

% ── Estructura ────────────────────────────────────────────────────────────────
\begin{frame}{Estructura de un proyecto de diapositivas}
  \begin{columns}[T]
    \begin{column}{0.45\textwidth}
      \footnotesize
      \texttt{slides/congreso/}\\
      \texttt{├── main.tex}\\
      \texttt{├── references.bib}\\
      \texttt{├── figures/}\\
      \texttt{└── slides/}\\
      \texttt{\phantom{└── }├── 00\_title.tex}\\
      \texttt{\phantom{└── }├── 01\_intro.tex}\\
      \texttt{\phantom{└── }└── 02\_content.tex}
    \end{column}
    \begin{column}{0.52\textwidth}
      \begin{itemize}
        \item \textbf{main.tex} — metadatos y orden de slides
        \item \textbf{slides/} — un archivo por seccion
        \item \textbf{figures/} — imagenes y graficos
        \item \textbf{references.bib} — bibliografia
      \end{itemize}
      \vspace{0.5em}
      \begin{block}{Tip}
        Agregar una seccion = crear un archivo y un \texttt{\textbackslash input\{\}} en \texttt{main.tex}.
      \end{block}
    \end{column}
  \end{columns}
\end{frame}

% ── Motores ───────────────────────────────────────────────────────────────────
\begin{frame}{Motores disponibles}
  \begin{table}
    \centering
    \begin{tabular}{llp{5cm}}
      \toprule
      Motor    & Flag                         & Cuando usarlo \\
      \midrule
      XeLaTeX  & (default)                    & Siempre. Unicode y fuentes del sistema \\
      LuaLaTeX & \texttt{-{}-engine lualatex} & Programacion avanzada con Lua \\
      pdfLaTeX & \texttt{-{}-engine pdflatex} & Documentos legacy \\
      \bottomrule
    \end{tabular}
  \end{table}

  \vspace{0.5em}
  \alert{XeLaTeX es el recomendado} para todos los templates de pptex.
\end{frame}
